% -------------------------------------------------------------------------
% Chapter 4: Watermark Evaluation and Robustness Analysis
% -------------------------------------------------------------------------

\chapter{Watermark Evaluation and Robustness Analysis}
\label{chap:watermark-results}

This chapter reports the empirical evaluation of the proposed \emph{zero-bit}
speech watermarking module introduced in Chapter~3. The goal of the module is
to embed an imperceptible provenance signal into assistant-generated TTS speech
and to enable fast, key-based verification of \emph{assistant-originated} audio.
Following the experimental protocol defined in Section~3.6, we focus on two
questions:

\begin{itemize}
    \item \textbf{RQ1 (Imperceptibility):} Does watermark embedding noticeably
    degrade the perceived quality of TTS speech?
    \item \textbf{RQ2 (Robustness):} Does the watermark remain detectable after
    realistic signal and channel distortions, including compression and
    re-recording-like effects?
\end{itemize}

The results in this chapter are obtained using the benchmarking script
\texttt{bench\_wm.py}, which systematically applies a suite of attacks to both
clean and watermarked audio and evaluates watermark presence using a fixed
decision threshold from the watermark configuration file.

\section{How the Proposed Watermark Works}
\label{sec:wm-how-it-works}

This section provides an implementation-oriented description of the watermark
scheme, complementing the high-level design overview in Section~3.3. The
watermark is \emph{zero-bit}: instead of transmitting a payload, the system only
needs to decide whether an audio clip contains the assistant's secret watermark
(\emph{present}) or not (\emph{absent}). This design is appropriate for provenance
checks in assistant pipelines, where the primary question is whether a signal was
produced by a trusted TTS component.

\subsection{Embedding as spread-spectrum energy modulation}
\label{subsec:wm-embed}

Let $x(t)$ be a mono waveform produced by the assistant's TTS model. The
embedder performs the following steps:

\begin{enumerate}
    \item \textbf{Resampling and framing.} The waveform is resampled to a fixed
    sample rate $f_s$ (in our experiments, $f_s=16$~kHz) and split into
    overlapping frames. A short-time Fourier transform (STFT) is computed to
    obtain a time--frequency representation.
    \item \textbf{Keyed bin selection.} A secret key $K$ deterministically
    defines a pseudo-random selection of time--frequency locations
    $S \subseteq \{(m,k)\}$ (frame index $m$, frequency bin $k$) that will carry
    watermark energy. Intuitively, $S$ specifies \emph{where} the watermark
    ``lives'' in the spectrogram.
    \item \textbf{Low-amplitude perturbation.} For the selected locations in
    $S$, the embedder applies a small perturbation to the spectral magnitudes.
    In a generic spread-spectrum formulation, this can be expressed as
    \[
        |X'(m,k)| = |X(m,k)| \cdot (1 + \alpha \, w_{m,k}), \quad (m,k)\in S,
    \]
    where $w_{m,k}\in\{-1,+1\}$ is a pseudo-random sign pattern derived from
    $K$, and $\alpha$ is a small embedding strength chosen to preserve
    imperceptibility. Locations outside $S$ are left unchanged.
    \item \textbf{Inverse transform.} The modified spectrum is converted back
    to the time domain via inverse STFT and overlap--add, yielding a watermarked
    waveform $x_{\mathrm{wm}}(t)$.
\end{enumerate}

Because energy is distributed across many time--frequency bins, the watermark
is \emph{spread-spectrum}: each individual modification is weak and typically
inaudible, but evidence can be aggregated over time at detection. In principle,
longer utterances provide more watermark observations and can improve
detectability by reducing variance of the detection statistic.

\subsection{Detection as a correlation-like statistic}
\label{subsec:wm-detect}

Given a candidate waveform $\tilde{x}(t)$, the verifier repeats the same
pre-processing steps (resampling, STFT computation, and keyed selection of
locations $S$) and computes a scalar detection score $s(\tilde{x})$. The score is
designed to be high if the expected keyed pattern is present and low otherwise.
A typical statistic in spread-spectrum watermarking is a normalized correlation:

\[
s(\tilde{x}) \propto \frac{1}{|S|}\sum_{(m,k)\in S} w_{m,k}\,\phi(\tilde{X}(m,k)),
\]

where $\phi(\cdot)$ denotes the chosen feature transform (e.g., magnitude,
log-magnitude, or a locally normalized residual). The exact form in our
implementation is correlation-based and yields a continuous score.

The verifier compares the score to a fixed threshold $\tau$:

\[
\text{watermark present} \iff s(\tilde{x}) \ge \tau.
\]

A high threshold reduces false alarms (clean audio incorrectly classified as
watermarked), while a low threshold improves true detection at the cost of more
false alarms. Section~\ref{sec:wm-robustness-results} reports all results at the
fixed threshold used in the benchmark configuration.

\subsection{Security assumptions}
\label{subsec:wm-assumptions}

Consistent with the threat model in Chapter~1, the watermark is keyed. We assume
attackers do not have white-box access to $K$ or to the exact embedding and
detection code. Under this assumption, an adversary attempting to impersonate
the assistant must either: (i) obtain the key, or (ii) remove the watermark using
generic signal processing operations (filtering, re-encoding, denoising, etc.),
which typically also degrade audio quality. The experiments below evaluate
precisely these generic transformations.

\section{Evaluation Setup}
\label{sec:wm-eval-setup}

\subsection{Dataset}
\label{subsec:wm-dataset}

The robustness benchmark was run on a set of $N=1001$ assistant-style
utterances (mono WAV files) generated by the TTS pipeline and resampled to the
watermark sample rate. For each utterance, the benchmark constructs both a
\emph{clean} version $x(t)$ and a \emph{watermarked} version
$x_{\mathrm{wm}}(t)$, yielding $2N$ base signals.

\subsection{Attack suite}
\label{subsec:wm-attacks}

Each base signal is subjected to a collection of common distortions intended to
approximate real-world channels and naive watermark removal attempts. The suite
covers:

\begin{itemize}
    \item \textbf{Identity (no attack).}
    \item \textbf{Amplitude transforms:} gain changes ($\pm 6$~dB) and hard
    clipping (peak limits 0.8 and 0.6).
    \item \textbf{Bandwidth and sampling effects:} resampling round-trips
    (downsample then upsample back) to 16~kHz and 8~kHz; low-pass filtering
    with cutoffs 7~kHz and 3.4~kHz (the latter approximates telephony
    bandwidth).
    \item \textbf{Additive noise:} white noise at several SNR levels
    (30/20/10/5~dB).
    \item \textbf{Reverberation-like effects:} convolution with a synthetic
    impulse response (exponential decay with random early reflections) at
    $RT_{60}\in\{0.2,0.4\}$ seconds.
    \item \textbf{Lossy compression:} codec round-trips via \texttt{ffmpeg} for
    MP3 (320/128/64/32~kbps) and AAC (128/64~kbps).
\end{itemize}

Overall, the benchmark evaluates 21 attack conditions per signal. Because both
clean and watermarked versions are evaluated, each attack yields
$2N=2002$ scored trials.

\subsection{Metrics}
\label{subsec:wm-metrics}

For each trial, the detector produces a scalar score $s$ and a binary decision at
the configured threshold $\tau$. We report (TPR/FPR are shown as fractions in $[0,1]$, i.e., multiply by 100 for percent):

\begin{itemize}
    \item \textbf{True Positive Rate (TPR):} fraction of watermarked signals
    classified as watermarked:
    $\mathrm{TPR}=\Pr[s\ge\tau \mid \text{watermarked}]$.
    \item \textbf{False Positive Rate (FPR):} fraction of clean signals
    incorrectly classified as watermarked:
    $\mathrm{FPR}=\Pr[s\ge\tau \mid \text{clean}]$.
    \item \textbf{Mean score shift:} mean detector score for watermarked and
    clean audio under each attack, $\overline{s}_{wm}$ and $\overline{s}_{clean}$,
    which provides additional insight even when TPR/FPR at a fixed threshold are
    low.
\end{itemize}

\section{Imperceptibility Evaluation (RQ1)}
\label{sec:wm-imperceptibility}

Robustness alone is insufficient for watermark deployment: assistant responses
must remain natural and free of audible artifacts. The thesis methodology
(Section~3.6.1) proposes two complementary subjective protocols:

\begin{itemize}
    \item \textbf{Mean Opinion Score (MOS):} listeners rate individual utterances
    on a five-point quality scale.
    \item \textbf{A/B preference:} listeners compare paired clean vs. watermarked
    utterances and select the preferred sample or ``no preference''.
\end{itemize}

At the time of writing, the benchmark results available for this chapter focus on
\emph{robustness} (RQ2) and do not yet include a completed human-subject
listening study. Therefore, the chapter does not claim a definitive answer to
RQ1. In the final thesis version, this section should be extended with:

\begin{enumerate}
    \item the number of participants and listening conditions (headphones,
    quiet room, online platform, etc.);
    \item MOS statistics for clean vs. watermarked speech (mean $\pm$ standard
    deviation, significance tests); and
    \item A/B preference outcomes (fraction preferring clean, preferring
    watermarked, and no preference).
\end{enumerate}

Nevertheless, the robustness results in Section~\ref{sec:wm-robustness-results}
are informative for the imperceptibility--robustness trade-off: the low baseline
TPR at a conservative threshold suggests that watermark strength may currently be
tuned cautiously, potentially prioritizing imperceptibility at the expense of
detectability.

\section{Robustness Results}
\label{sec:wm-robustness-results}

The benchmark used a fixed watermark decision threshold of $\tau=0.02$ (as
configured in \texttt{watermark.yaml}). Table~\ref{tab:wm-robustness-full}
summarizes detection performance for each attack.

\begin{table}[t]
\centering
\small
\caption{Zero-bit watermark detection results at fixed threshold $\tau=0.02$
on $N=1001$ utterances. For each attack we report TPR/FPR and mean detector
scores for watermarked and clean signals. Because each attack is evaluated on
both clean and watermarked versions, each condition has $n=2002$ trials.}
\label{tab:wm-robustness-full}
\begin{tabular}{l l r r r r}
\hline
\textbf{Attack} & \textbf{Params} &
\textbf{TPR} & \textbf{FPR} &
$\boldsymbol{\overline{s}_{wm}}$ & $\boldsymbol{\overline{s}_{clean}}$ \\
\hline
none & -- & 0.0609 & 0.0000 & 0.0133 & -0.0286 \\
gain & $+6$ dB & 0.0529 & 0.0000 & 0.0130 & -0.0286 \\
gain & $-6$ dB & 0.0659 & 0.0000 & 0.0135 & -0.0286 \\
clip & 0.8 & 0.0609 & 0.0000 & 0.0133 & -0.0286 \\
clip & 0.6 & 0.0589 & 0.0000 & 0.0133 & -0.0286 \\
resample\_roundtrip & 16 kHz & 0.0609 & 0.0000 & 0.0133 & -0.0286 \\
resample\_roundtrip & 8 kHz & 0.0000 & 0.0000 & -0.0467 & -0.0479 \\
lowpass & 7000 Hz & 0.0160 & 0.0000 & 0.0113 & -0.0295 \\
lowpass & 3400 Hz & 0.0000 & 0.0000 & -0.0418 & -0.0460 \\
white\_noise & 30 dB SNR & 0.0000 & 0.0000 & $1.9\times 10^{-5}$ & -0.0203 \\
white\_noise & 20 dB SNR & 0.0000 & 0.0000 & 0.0022 & -0.0135 \\
white\_noise & 10 dB SNR & 0.0000 & 0.0000 & 0.0039 & -0.0057 \\
white\_noise & 5 dB SNR & 0.0000 & 0.0000 & 0.0040 & -0.0026 \\
reverb\_synth & $RT_{60}=0.2$ & 0.0000 & 0.0000 & -0.0041 & -0.0475 \\
reverb\_synth & $RT_{60}=0.4$ & 0.0000 & 0.0000 & -0.0050 & -0.0466 \\
mp3\_roundtrip & 320 kbps & 0.2837 & 0.0000 & 0.0178 & -0.0274 \\
mp3\_roundtrip & 128 kbps & 0.2997 & 0.0000 & 0.0179 & -0.0274 \\
mp3\_roundtrip & 64 kbps & 0.1708 & 0.0000 & 0.0162 & -0.0271 \\
mp3\_roundtrip & 32 kbps & 0.0020 & 0.0000 & 0.0086 & -0.0269 \\
aac\_roundtrip & 128 kbps & 0.0969 & 0.0000 & 0.0147 & -0.0271 \\
aac\_roundtrip & 64 kbps & 0.0909 & 0.0000 & 0.0145 & -0.0271 \\
\hline
\end{tabular}
\end{table}

\subsection{Baseline performance and threshold behavior}
\label{subsec:wm-baseline}

In the \emph{no-attack} condition, the watermark achieves a TPR of
$0.0609$ (61 detections out of 1001 watermarked utterances), while maintaining
an FPR of $0$ (0 false alarms out of 1001 clean utterances). This indicates that
the current operating point is extremely conservative: clean audio is strongly
separated from the threshold (mean score $\overline{s}_{clean}\approx -0.0286$),
but the average watermark score $\overline{s}_{wm}\approx 0.0133$ remains below
$\tau=0.02$. As a result, only the upper tail of the watermarked score
distribution exceeds the threshold.

Operationally, this highlights an important trade-off: the system currently favors
\emph{no false alarms} over \emph{high recall}. For a provenance check, false
alarms are indeed undesirable (they could incorrectly authenticate attacker
audio), but false negatives can harm usability if legitimate assistant audio is
frequently rejected. A more complete characterization would include ROC/DET
analysis to select $\tau$ based on a target operating point (e.g., a specified
upper bound on FPR). In the present evaluation we keep $\tau$ fixed to isolate
the effect of channel distortions.

\subsection{Robustness to simple amplitude distortions}
\label{subsec:wm-amplitude}

Gain changes of $\pm 6$~dB produce only small changes in TPR (5.29\% and
6.59\%, respectively), and clipping at 0.8/0.6 similarly leaves baseline
performance nearly unchanged. This behavior is consistent with a detector that
relies on relative spectral patterns or internally normalized statistics rather than
absolute waveform amplitude.

\subsection{Vulnerability to bandwidth reduction and resampling}
\label{subsec:wm-bandwidth}

Bandwidth reduction is a critical failure mode. Resampling through an 8~kHz
bottleneck (downsample to 8~kHz, then upsample back) yields a TPR of 0 at the
fixed threshold. Likewise, a 3.4~kHz low-pass filter (telephony-like bandwidth)
also reduces TPR to 0.

These results strongly suggest that, in its current configuration, the watermark's
detectable energy is concentrated in frequency regions that are removed or
distorted by narrowband channels. In practical terms, an attacker can remove the
watermark using a \emph{low-effort} transformation (downsampling or aggressive
low-pass filtering) while still preserving intelligibility of the speech content.
This vulnerability directly impacts deployment scenarios where assistant audio
passes through bandwidth-limited transmission paths or is re-recorded through
devices with limited high-frequency response.

A milder low-pass filter at 7~kHz (still within the 16~kHz sampling regime)
reduces TPR from 6.09\% to 1.60\%, indicating that even partial attenuation of
the upper band can substantially weaken detection.

\subsection{Vulnerability to additive noise}
\label{subsec:wm-noise}

Under additive white noise, the watermark is not detected at all at the fixed
threshold for SNR values from 30~dB down to 5~dB (TPR=0 across all tested
levels). While the mean watermark score remains slightly positive at lower SNR
levels (e.g., $\overline{s}_{wm}\approx 0.004$ at 5~dB), it stays well below the
configured threshold.

Notably, increasing noise moves the \emph{clean} mean score toward zero
($\overline{s}_{clean}$ increases from approximately $-0.0286$ to $-0.0026$ at
5~dB). This indicates a potential risk when lowering the threshold to improve
TPR: noise may simultaneously increase false alarms by shifting clean scores
upward. This interaction suggests that robust watermarking in noisy environments
requires either stronger embedding, more sophisticated normalization, or
augmentation-aware thresholding.

\subsection{Vulnerability to reverberation and convolutional effects}
\label{subsec:wm-reverb}

Synthetic reverberation (modeled as convolution with a random impulse response
with $RT_{60}\in\{0.2,0.4\}$) yields TPR=0 at the fixed threshold. Unlike the
noise case, the watermarked mean score remains higher than the clean mean score
(e.g., $-0.0041$ vs. $-0.0475$ at $RT_{60}=0.2$), indicating that some watermark
evidence may still be present but is substantially weakened and shifted below the
decision boundary.

This pattern is consistent with time--frequency smearing: reverberation introduces
delayed replicas of the signal, altering frame-local spectral structure and
reducing the correlation with the keyed template.

\subsection{Behavior under lossy audio codecs}
\label{subsec:wm-codecs}

Lossy compression shows mixed behavior. MP3 round-trips at 128--320~kbps yield
the strongest results in this benchmark (TPR $\approx 0.28$--$0.30$), while
lower MP3 bitrates degrade sharply (TPR 0.17 at 64~kbps and near zero at
32~kbps). AAC yields moderate TPR values around 9--10\% at 64--128~kbps.

Interestingly, for high-bitrate MP3 the measured TPR is \emph{higher} than in the
no-attack case. The mean watermark score also increases (from 0.0133 to
$\approx 0.0179$), suggesting that the codec round-trip can shift the score
distribution upward, making more watermarked samples cross the fixed threshold.
One plausible explanation is that codec quantization and psychoacoustic shaping
alter spectral statistics in a way that better aligns with the detector's feature
normalization. However, despite this improvement, the recall remains below 31\%,
and therefore the watermark is still unreliable as a standalone provenance signal
under the current operating point.

\section{Summary of Current Vulnerabilities}
\label{sec:wm-vuln-summary}

The robustness benchmark supports the following conclusions about the current
implementation:

\begin{itemize}
    \item \textbf{Strong at avoiding false alarms (at $\tau=0.02$).} FPR is 0 for
    all tested attacks, meaning no clean samples exceed the threshold. This is a
    desirable property for provenance verification when the cost of falsely
    authenticating attacker audio is high.
    \item \textbf{Weak true detection even without channel distortion.} Baseline
    TPR is only 6.09\%. This indicates that the watermark is either embedded very
    conservatively (to preserve imperceptibility), the threshold is set too high,
    or implementation details reduce effective watermark strength.
    \item \textbf{Highly vulnerable to bandwidth limitation.} Resampling through
    8~kHz and low-pass filtering at 3.4~kHz fully remove detectability at the
    fixed threshold. This implies that the watermark relies on spectral regions
    that are absent in narrowband channels.
    \item \textbf{Highly vulnerable to additive noise and reverberation.} In the
    tested regime, both noise and reverberation lead to TPR=0 at the fixed
    threshold, which is problematic for realistic far-field and re-recording
    scenarios.
    \item \textbf{Partially robust to moderate codec compression.} MP3 and AAC
    at moderate to high bitrates preserve some detectability, but robustness
    degrades at low bitrates.
\end{itemize}

Overall, these findings suggest that the current watermark configuration is
\emph{not yet robust} to realistic channel degradations expected in voice
assistant deployments, especially when the audio may be transmitted over
bandwidth-limited links or captured in noisy and reverberant environments.

\section{Chapter Conclusion}
\label{sec:wm-chapter-conclusion}

This chapter evaluated the proposed zero-bit spread-spectrum watermark under a
broad suite of signal and channel distortions. At the current threshold
$\tau=0.02$, the watermark exhibits \emph{zero false alarms} across all tested
conditions, but also exhibits \emph{low true detection} in clean conditions and
fails under several realistic degradations (narrowband resampling, low-pass
filtering, additive noise, and reverberation). These results identify clear
failure modes and motivate subsequent work on (i) threshold selection via ROC/DET
analysis, (ii) multi-band or bandwidth-aware embedding, and (iii) improved
normalization and augmentation strategies to maintain detectability under
noise and re-recording.

\section*{References}
\small
\noindent [1] [15] (PDF) On the Security and Privacy Challenges of Virtual Assistants. \url{https://www.researchgate.net/publication/350413576_On_the_Security_and_Privacy_Challenges_of_Virtual_Assistants}\par
\noindent [2] [6] [7] [8] [13] [17] [18] [38] [39] [40] [41] [42] [43] [44] [45] [48] [2308.14970] Audio Deepfake Detection: A Survey. \url{https://ar5iv.labs.arxiv.org/html/2308.14970}\par
\noindent [3] [16] usenix.org. \url{https://www.usenix.org/system/files/usenixsecurity24-jiang-nan.pdf}\par
\noindent [4] [5] [12] [23] [24] [29] [30] [31] [49] usenix.org. \url{https://www.usenix.org/system/files/usenixsecurity25-zong.pdf}\par
\noindent [9] [11] [19] [20] [21] [22] Audio Watermarking Won’t Solve the Real Dangers of AI Voice Manipulation -- Center for Data Innovation. \url{https://datainnovation.org/2024/10/audio-watermarking-wont-solve-the-real-dangers-of-ai-voice-manipulation/}\par
\noindent [10] [28] [35] [36] [37] [46] [47] [50] [51] [52] [53] [54] FakeMark: Deepfake Speech Attribution With Watermarked Artifacts. \url{https://arxiv.org/html/2510.12042v2}\par
\noindent [14] Proactive Detection of Voice Cloning with Localized Watermarking. \url{https://ai.meta.com/research/publications/proactive-detection-of-voice-cloning-with-localized-watermarking/}\par
\noindent [25] [26] [27] [2401.17264] Proactive Detection of Voice Cloning with Localized Watermarking. \url{https://arxiv.org/abs/2401.17264}\par
\noindent [32] [33] [34] [2504.14832] Protecting Your Voice: Temporal-aware Robust Watermarking. \url{https://arxiv.org/abs/2504.14832}\par